% Bagian: computability
% Bab: computability-theory
% Subbagian: k-1

\documentclass[../../../include/open-logic-section]{subfiles}

\begin{document}

\olfileid[id]{cmp}{thy}{k1}
\olsection{Contoh Redusibilitas}

Mari kita tinjau suatu penerapan \olref[ppr]{prop:reduce}.

\begin{prop}
\ollabel{prop:k1}
Misalkan
\[
K_1 = \Setabs{e}{\cfind{e}(0) \fdefined}.
\]
Maka $K_1$ terenumerasi secara komputabel, tetapi tidak dapat dihitung.
\end{prop}

\begin{proof}
Karena $K_1 = \Setabs{e}{\lexists[s][T(e,0,s)]}$, $K_1$ terenumerasi
secara komputabel menurut \olref[eqc]{thm:exists-char}.

Untuk menunjukkan bahwa $K_1$ tidak dapat dihitung, mari kita tunjukkan bahwa
$K_0$ tereduksi kepadanya.

\begin{explain}
Hal ini agak rumit karena dengan menggunakan $K_1$ kita hanya dapat mengajukan
pertanyaan tentang komputasi yang dimulai dengan masukan tertentu, yaitu~$0$.
Andaikan Anda mempunyai teman cerdas yang dapat menjawab pertanyaan semacam
ini (teman seperti itu dikenal sebagai ``orakel''). Lalu, andaikan seseorang
mendatangi Anda dan bertanya apakah $\tuple{e,x}$ berada dalam $K_0$, yakni
apakah mesin~$e$ berhenti pada masukan~$x$. Salah satu hal yang dapat Anda
lakukan ialah membangun mesin lain, $e_x$, yang untuk \emph{sembarang} masukan
mengabaikan masukan tersebut dan sebagai gantinya menjalankan~$e$ pada
masukan~$x$. Maka jelas bahwa pertanyaan apakah mesin~$e$ berhenti pada
masukan~$x$ ekuivalen dengan pertanyaan apakah mesin~$e_x$ berhenti pada
masukan~$0$ (atau masukan lain apa pun). Jadi, Anda bertanya kepada teman Anda
apakah mesin baru~$e_x$ ini berhenti pada masukan~$0$; jawaban teman Anda atas
pertanyaan yang diubah itu memberikan jawaban atas pertanyaan semula. Ini
memberikan reduksi yang diinginkan dari $K_0$ ke~$K_1$.
\end{explain}

Dengan menggunakan fungsi dapat dihitung parsial universal, misalkan $f$
fungsi berargumen~3 yang didefinisikan oleh
\[
f(x,y,z) \simeq \cfind{x}(y).
\]
Perhatikan bahwa $f$ sama sekali mengabaikan masukan ketiganya. Pilih suatu
indeks $e$ sedemikian sehingga $f = \cfind{e}[3]$; jadi kita mempunyai
\[
\cfind{e}[3](x,y,z) \simeq \cfind{x}(y).
\]
Menurut teorema $s$-$m$-$n$, terdapat fungsi $s(e,x,y)$ sedemikian sehingga,
bagi setiap $z$,
\begin{align*}
\cfind{s(e,x,y)}(z) & \simeq \cfind{e}[3](x,y,z) \\
& \simeq \cfind{x}(y).
\end{align*}

\begin{explain}
Dalam argumen informal di atas, $s(e,x,y)$ merupakan indeks bagi mesin yang,
untuk sembarang masukan~$z$, mengabaikan masukan itu dan menghitung
$\cfind{x}(y)$.
\end{explain}

Khususnya, kita mempunyai
\[
\cfind{s(e,x,y)}(0) \fdefined \quad \text{jika dan hanya jika} \quad
\cfind{x}(y) \fdefined.
\]
Dengan kata lain, $\tuple{x,y} \in K_0$ jika dan hanya jika $s(e,x,y) \in
K_1$. Jadi, fungsi $g$ yang didefinisikan oleh
\[
g(w) = s(e,(w)_0,(w)_1)
\]
merupakan reduksi dari $K_0$ ke~$K_1$.
\end{proof}

\end{document}
